These notes have two chapters.
Chapter \ref{chap.microstructure} is about market microstructure:
what an order is,
how an exchange matches orders,
what the limit order book is at any instant,
and how the flow of orders forms a price.
It is the subject of the course.
Chapter \ref{chap.options} is option pricing and volatility trading:
the Black--Scholes model,
the fundamental theorem of derivative trading,
and the SABR model.
It is optional reading for the student who wants to go further:
it is neither taught in class nor examined,
and it is added to these notes towards the end of the course.

The course is taught in Python.
Every statement of these notes that can be computed is computed,
and students are expected to read and understand all the implementations that compute them.
The code is available in the course repository,
\url{https://github.com/claudio-ICL/unito26}.

The first half of the course presents Chapter \ref{chap.microstructure} in lectures.
The second half is laboratory work on its implementation:
we take a closer look at the code, and we read it, run it and change it together.
Students are strongly encouraged to write to the lecturer
with questions and remarks on what they find challenging in the implementation,
and the laboratory sessions concentrate on those parts.
The implementation is organised in classes, which the course does not assume:
Appendix \ref{chap.objectOrientedProgramming} introduces them on the code of the repository.

The angle is the one the practice has moved to.
Writing code that works was, until recently, the skill that gated everything else.
Nowadays, with the advent of coding assistants,
writing functioning code no longer gates quantitative analysis.
However,
coding assistants do not supply the \emph{judgment}
to tell a correct implementation from one that merely runs,
to choose between two designs,
or
to understand the long-term impact of an implementation choice.
The course teaches this judgment,
and the exam will evaluate the students on that.
